\section{Conclusion}
\label{sec:conclusion}

We introduced AreaSection, a dual-constraint extension of minimum-edge-cut
redistricting that simultaneously balances population and land area in each
bisection.
The algorithm uses METIS's native $\mathtt{ncon}=2$ multi-constraint
partitioning with isoperimetric normalisation for ratio selection and a
Lorenz-curve pre-filter for feasibility.

The empirical results across 50 states show a consistent pattern:
\textbf{AreaSection dramatically reduces urban peeling for most states}
(Illinois: 1:16 $\to$ 8:9; New York: 1:25 $\to$ 13:13;
Florida: 1:27 $\to$ 14:14) while retaining asymmetric splits for states
where metropolitan concentration makes equal geographic splits infeasible
(Georgia: 3:11, Massachusetts: 2:7).

\textbf{The partisan outcome is mostly stable}: in the 34-state head-to-head
comparison, 25 states produce identical seat allocations and 9 differ by one
seat.
Wisconsin is stable under both algorithms; North Carolina shifts from
5D/9R under GeoSection to 6D/8R under AreaSection.
AreaSection is a geometric regularity mechanism, not a partisan mechanism.
It changes the \emph{competitiveness profile} of districts and, in a minority
of states, can change a marginal seat, but the direction is not systematic.

This stability finding supports the broader claim of the Algorithmic
Redistricting project: for states where population distribution determines
the partisan outcome, the specific bisection algorithm is a second-order
choice.
The first-order question is: does the state's geography allow enough
competitive districts for the minority party to win seats proportional to
its vote share?
For Wisconsin and many other states, the answer is no regardless of whether
the algorithm is GeoSection or AreaSection.
North Carolina shows the boundary case: a different first split can recover
one Democratic seat, but not enough to eliminate the underlying geographic
proportionality gap.

\paragraph{Future work.}
Three directions merit further investigation:

\begin{itemize}
\item \textbf{Recursive area constraints.} We apply the area constraint only
  at the first bisection level. Applying it at every level would enforce
  geographic balance throughout the bisection tree.
  Preliminary results suggest this improves balance in large states but
  increases the balance tolerance failures.

\item \textbf{Area swing parameter sensitivity.} Systematic testing of
  different $\mathtt{ubvec}[1]$ values (from $\pm 5\,\%$ to $\pm 25\,\%$)
  to characterise the tradeoff between geographic balance and population
  balance feasibility.

\item \textbf{Multi-year comparison.} Applying AreaSection to 2000 and 2010
  Census data to test stability of the natural ratio structure across
  census cycles, as population distribution shifts due to urbanisation trends.
\end{itemize}

\paragraph{Code and data availability.}
The AreaSection implementation is part of the open-source \texttt{bisect}
Rust binary, available at \url{https://github.com/giodl73-repo/REDIST}.
All 50-state results and tract-level assignments are available in the
replication archive (see \texttt{research/A.4+replication-materials}).
